\documentclass[11pt,a4paper]{article}
\usepackage[margin=1in]{geometry}
\usepackage{hyperref}
\usepackage{enumitem}
\title{Strategic Management \\ Week 10 \\ Slide-by-Slide Notes (ELI5) \\ Unit 6}
\author{(Generated from \texttt{sm-week-10-slides-6.pdf})}
\date{}

\begin{document}
\maketitle

\section*{How to use this (ELI5)}
\begin{itemize}
  \item For every slide, write the ``Exam write'' line in your notebook.
  \item Then copy the ELI5 explanation in simple English.
  \item Use it for long answers: label + reasons + risks + recommendation logic.
\end{itemize}

% ---------------------------
\section*{Slide 1 of 13 (Contents: Unit 6)}
\textbf{Key idea (from slide):} Unit 6 is Corporate Strategy Development (where to compete).\\
\textbf{ELI5:} This part is about choosing the company map: what businesses you want.\\
\textbf{Exam write:} ``Unit 6 answers where to compete (scope, growth moves, and restructuring).''

% ---------------------------
\section*{Slide 2 of 13 (Learning objectives)}
\textbf{Key idea (from slide):}
\begin{itemize}
  \item key concepts of corporate strategy: development and methods
  \item answer: where to compete
  \item analyze development strategies: expansion, diversification, vertical integration
  \item reasons and ways of restructuring businesses
\end{itemize}
\textbf{ELI5:} You learn the big corporate moves and how to fix poor businesses.\\
\textbf{Exam write:} ``I will match corporate-scope choices to reasons and restructuring alternatives.''

% ---------------------------
\section*{Slide 3 of 13 (Firm scope, development and growth)}
\textbf{Key idea (from slide):}
\begin{itemize}
  \item firm development = quantitative and qualitative changes from scope changes
  \item corporate strategies include specialization, diversification (related/unrelated), vertical integration, restructuring, internationalization
  \item firm growth = increases/decreases in assets, outputs, sales, profitability, etc.
  \item corporate methods: internal growth, M\&A, strategic alliances
\end{itemize}
\textbf{ELI5:} Corporate strategy says: ``What areas will we play in?'' and ``How will we grow there?''\\
\textbf{Exam write:} ``Corporate development changes scope; corporate growth changes size variables.''

% ---------------------------
\section*{Slide 4 of 13 (Firm scope addresses where to compete)}
\textbf{Key idea (from slide):} Firm scope = set of businesses the company wants to compete in; key question: WHERE TO COMPETE?\\
\textbf{ELI5:} Scope is like your playlist: which industries you focus on.\\
\textbf{Exam write:} ``In scope questions, define the set of businesses and answer where the firm competes.''

% ---------------------------
\section*{Slide 5 of 13 (Specialization title)}
\textbf{Key idea (from slide):} The slide starts specialization discussion.\\
\textbf{ELI5:} Specialization means you do one main thing better.\\
\textbf{Exam write:} ``Specialization focuses on the core where the firm is strongest.''

% ---------------------------
\section*{Slide 6 of 13 (Specialization reasons)}
\textbf{Key idea (from slide):} Reasons for specialization:
\begin{itemize}
  \item economies of scale
  \item consolidate market position; lower costs or more value added
  \item increase sales using existing products, new customers, improvements
  \item low-risk: same products for same needs
\end{itemize}
\textbf{ELI5:} It is safer because you stay close to what you already know.\\
\textbf{Exam write:} ``Specialization is justified by scale and by improving cost/value in the same customer-need.''

% ---------------------------
\section*{Slide 7 of 13 (Diversification title)}
\textbf{Key idea (from slide):} Diversification starts.\\
\textbf{ELI5:} Diversification means you do more than one type of business.\\
\textbf{Exam write:} ``Diversification changes scope by entering more than one business area.''

% ---------------------------
\section*{Slide 8 of 13 (Diversification reasons and risks)}
\textbf{Key idea (from slide):}
\begin{itemize}
  \item Reasons: economies of scope (synergies), improve competitive position, enter attractive industries,
  reduce risk, obtain internal funding, exploit managerial capabilities
  \item risk: synergies require coordination/commitment, can reduce flexibility, diversification can be costly (entry barriers, high internal funding cost)
  \item metaphor: do not put all eggs in one basket
\end{itemize}
\textbf{ELI5:} Diversification lowers danger if one business is weak, but it can be hard to manage and coordinate.\\
\textbf{Exam write:} ``Use diversification when synergies/risk reduction dominate, and show the coordination cost and flexibility risk.''

% ---------------------------
\section*{Slide 9 of 13 (Vertical integration definition)}
\textbf{Key idea (from slide):} Vertically (de)integrated if it undertakes several activities in the production process.\\
\textbf{ELI5:} Vertical integration means you also control steps before/after your product.\\
\textbf{Exam write:} ``Vertical integration is controlling multiple production-process activities inside the firm.''

% ---------------------------
\section*{Slide 10 of 13 (Vertical integration reasons and risks)}
\textbf{Key idea (from slide):}
\begin{itemize}
  \item reasons: easy coordination/planning; reduce opportunism risk from external suppliers/clients; increase customer value via quality/services; protect information/knowledge; control the final market
  \item risks: more risk doing more activities for same final market; less flexibility; higher costs from administration, lack of experience, and low production scale in each activity
\end{itemize}
\textbf{ELI5:} You integrate to control quality and information, but it can cost more and make you less flexible.\\
\textbf{Exam write:} ``Support vertical integration with coordination/quality/info reasons, and address cost and inflexibility risks.''

% ---------------------------
\section*{Slide 11 of 13 (Restructuring definition + alternatives)}
\textbf{Key idea (from slide):}
\begin{itemize}
  \item restructuring = modify or reduce firm scope, with possible withdrawal (divestment) from at least one business
  \item alternatives: specialization, focusing on core businesses; business portfolio restructuring (withdraw and/or enter businesses); business restructuring to re-float a poor-performance business
\end{itemize}
\textbf{ELI5:} Restructuring is fixing a bad business by changing the company map: keep, cut, or replace parts.\\
\textbf{Exam write:} ``Restructuring means redefining scope and choosing divestment/entry or turnaround actions.''

% ---------------------------
\section*{Slide 12 of 13 (Restructuring: two basic options)}
\textbf{Key idea (from slide):}
\begin{itemize}
  \item re-floating: possible solutions/measures include change managers, redefine competitive strategy, sell assets, refinance debt, cost reduction
  \item withdrawal: no full solution; alternatives include sale, harvest (max short-term financial flows), winding-up (end operations and sell remaining assets)
\end{itemize}
\textbf{ELI5:} If turnaround is possible, you re-float; if not, you sell/harvest/close.\\
\textbf{Exam write:} ``In poor performance cases, choose turnaround measures vs withdrawal based on whether recovery is feasible.''

% ---------------------------
\section*{Slide 13 of 13 (Why restructuring happens: causes)}
\textbf{Key idea (from slide):}
\begin{itemize}
  \item when a business does not meet profit expectations: poor performance (sales/profits down), heavy borrowing
  \item causes: inefficient management, new competitors, unexpected demand change, bad competitive strategy, high internal costs, organizational inertia, etc.
\end{itemize}
\textbf{ELI5:} Restructuring is triggered because performance is bad and reasons must be diagnosed.\\
\textbf{Exam write:} ``Explain causes of poor performance to justify whether turnaround (re-float) or withdrawal is appropriate.''

\end{document}

